\newpage
\begin{center}
	\textbf{\large 3. Техническое описание}
\end{center}
\refstepcounter{chapter}
\addcontentsline{toc}{chapter}{3. Техническое описание}

Задача локализации QR-кода возникает в контексте инвентаризации оборудования в дата-центре: на инвентарной метке размещён QR-код, который требуется автоматически обнаружить на фотографии стойки или фронтальной панели устройства.
Особенностью данных является большое разрешение исходных изображений при малом относительном размере QR-кода.
Цель алгоритма -- определить границы QR-кода в координатах исходного изображения и представить результат в виде четырёх вершин многоугольника, аппроксимирующего область QR-кода без включения \textit{quiet zone}.

\paragraph{Основные определения и обозначения.}
Пусть входное изображение задано как матрица яркости $I(x, y)$ (для цветного изображения предварительно выполняется преобразование в оттенки серого).
Требуется определить четыре точки
\[
	Q = \{q_{TL};\ q_{TR};\ q_{BR};\ q_{BL}\},
\]
где $q_{TL}=(x_{TL}, y_{TL})$ -- верхняя левая вершина, далее вершины перечисляются по часовой стрелке.

Для оценки качества локализации применяются:
\begin{itemize}
	\item IoU (Intersection over Union) -- отношение площади пересечения двух многоугольников (предсказанного и разметки) к площади объединения;
	\item средняя ошибка по углам:
	      $$ e_{mean}=\frac{1}{4}\sum_{i=1}^{4} \|q_i - q_i^{GT}\| $$
	\item максимальная ошибка по углам:
	      $$ e_{max}=\max_{i\in\{1,2,3,4\}} \|q_i - q_i^{GT}\| $$
\end{itemize}

\paragraph{Идея метода: локализация по угловым маркерам.}
QR-код содержит три угловых маркера (finder patterns), расположенных в углах \texttt{TL}, \texttt{TR}, \texttt{BL}.
Каждый finder pattern имеет характерную вложенную структуру ``чёрный--белый--чёрный'' и форму, близкую к квадрату.
Использование finder patterns позволяет восстановить положение и масштаб QR-кода даже в случае, когда модули QR не образуют единую связную компоненту, а попытки выделить внешний контур напрямую приводят к ошибкам (например, выделяется только область finder).

Алгоритм состоит из двух основных этапов:
\begin{enumerate}
	\item поиск кандидатов на finder patterns;
	\item восстановление границ QR-кода по геометрии выбранной тройки finder patterns.
\end{enumerate}

\paragraph{Предобработка и уменьшение масштаба.}
Для снижения вычислительной сложности исходное изображение может быть уменьшено до максимального размера по стороне $M$ (например, $M\approx 1400$--$2200$ пикселей).
Уменьшение выполняется с интерполяцией \texttt{INTER\_AREA}, что даёт устойчивость при уменьшении и подавляет высокочастотный шум.

Далее выполняются:
\begin{itemize}
	\item сглаживание (Gaussian blur) для подавления шумов и облегчения бинаризации;
	\item бинаризация по Отсу (\texttt{THRESH\_OTSU}) для выделения контрастных областей.
\end{itemize}

\begin{figure}[!h]
	\begin{center}
		\includegraphics[width=250pt]{im09.jpg}
		\caption{Исходное изображение}
		\label{nmp10}
	\end{center}
\end{figure}

\begin{figure}[!h]
	\begin{center}
		\includegraphics[width=250pt]{im09-1.png}
		\caption{Grayscale изображение, сглаженное}
		\label{nmp11}
	\end{center}
\end{figure}

\begin{figure}[!h]
	\begin{center}
		\includegraphics[width=250pt]{im09-2.png}
		\caption{Бинаризированное изображение}
		\label{nmp12}
	\end{center}
\end{figure}

\paragraph{Поиск кандидатов finder patterns по контурам и иерархии.}
На бинарном изображении извлекаются контуры с построением дерева вложенности:
\[
	\texttt{findContours}(\cdot, \texttt{RETR\_TREE}, \texttt{CHAIN\_APPROX\_SIMPLE}).
\]
Finder pattern имеет вложенную структуру, поэтому в качестве кандидатов рассматриваются контуры, имеющие как минимум два уровня вложенности (наличие ``ребёнка'' и ``внука'' в иерархии контуров).

Для каждого такого контура выполняются геометрические проверки:
\begin{itemize}
	\item аппроксимация контура многоугольником \texttt{approxPolyDP};
	\item требование выпуклого четырёхугольника;
	\item проверка прямоугольности (углы близки к $90^\circ$);
	\item ограничение на отношение сторон (кандидат не должен быть вытянутым).
\end{itemize}

Для каждого кандидата вычисляется:
\begin{itemize}
	\item центр $c=(x_c,y_c)$;
	\item характерный размер $F$ (оценка стороны квадрата) по \texttt{minAreaRect}.
\end{itemize}

Далее выполняется дедупликация кандидатов по близости центров и ранжирование по площади.

\paragraph{Выбор тройки finder patterns и определение ролей TL/TR/BL.}
Из набора кандидатов перебираются тройки и выбирается та, которая наиболее соответствует геометрии QR-кода.
Критерий выбора основывается на двух свойствах:
\begin{enumerate}
	\item один из углов треугольника, построенного по центрам трёх finder patterns, должен быть близок к $90^\circ$ (это соответствует \texttt{TL});
	\item расстояния от \texttt{TL} до двух остальных маркеров должны быть сопоставимы (не должно быть сильной деградации масштаба в одном направлении).
\end{enumerate}

После выбора тройки вершина с углом, ближайшим к прямому, назначается как \texttt{TL}.
Два оставшихся маркера распределяются на \texttt{TR} и \texttt{BL} по ориентации в системе координат изображения (с учётом направления осей $x$ вправо и $y$ вниз).

\paragraph{Восстановление границ QR-кода по геометрии finder patterns.}
Поскольку finder pattern занимает $7\times7$ модулей QR-кода, расстояние от центра finder до внешней границы QR-кода по каждой оси составляет $3.5$ модуля.
Если оценка размера стороны finder pattern в пикселях равна $F$, то смещение от центра finder до внешнего края по нормализованным осям QR равно $F/2$.

Пусть $c_{TL}, c_{TR}, c_{BL}$ --- центры соответствующих finder patterns.
Определим единичные направления:
\[
	u = \frac{c_{TR}-c_{TL}}{\|c_{TR}-c_{TL}\|}, \quad
	v = \frac{c_{BL}-c_{TL}}{\|c_{BL}-c_{TL}\|}.
\]
Для повышения устойчивости при небольших искажениях выполняется ортогонализация $v$ относительно $u$.

Оценка стороны QR-кода (без \textit{quiet zone}):
\[
	S_x \approx \|c_{TR}-c_{TL}\| + F,\quad
	S_y \approx \|c_{BL}-c_{TL}\| + F,\quad
	S=\frac{S_x+S_y}{2}.
\]

Верхняя левая вершина внешней границы оценивается как:
\[
	q_{TL}=c_{TL}-\frac{F}{2}u-\frac{F}{2}v.
\]
Остальные вершины:
\[
	q_{TR}=q_{TL}+Su,\quad
	q_{BL}=q_{TL}+Sv,\quad
	q_{BR}=q_{TL}+Su+Sv.
\]
Полученный четырёхугольник упорядочивается в порядке \texttt{TL}, \texttt{TR}, \texttt{BR}, \texttt{BL}.

\paragraph{Вывод результата.}
Результат локализации сохраняется в JSON-файл, совместимый с форматом LabelMe: при успешной локализации в массив \texttt{shapes} добавляется один объект типа \texttt{polygon} с меткой \texttt{qr}, вершины которого соответствуют найденному четырёхугольнику.

\paragraph{Ограничения метода.}
Метод ориентирован на случаи, когда:
\begin{itemize}
	\item на изображении присутствует один QR-код;
	\item фон метки однородный, а QR-код имеет достаточный контраст;
	\item перспективные искажения невелики.
\end{itemize}
